\section{True operating updates from censored exits}

The actual operating clocks can change with the menu. We never infer them
from the clocks at balancing fields. Fix a requested coordinate, its
neighborhood, starting bit, and full operating environment $u$. Pin that
bit, release, and let its first-exit time be $T\sim\operatorname{Exp}(\lambda)$,
where $\lambda$ is the applicable rate in \eqref{eq:Q}. It obeys
$0<\lambda\le r^*$.

\begin{lemma}[Operating race and finite walk]\label{lem:walk}
Take $\rho=2r^*$, generate an independent
$E\sim\operatorname{Exp}(\rho)$, and censor the focal first-exit observation
at $E$. The bit $\one_{\{T<E\}}$ has probability
$p=\lambda/(\lambda+\rho)\le1/3$.
A random walk starting at 0, moving up 1 on a $p$-coin success and down 1
otherwise, ever hits $+1$ with probability
$p/(1-p)=\lambda/\rho$.
For $0<\theta\le1/4$, put
\[
 m=\lceil\log_2(8/\theta)\rceil,\qquad n=72m .
\]
Stop the walk at $+1$, at $-m$, or after $n$ steps; output 1 only if $+1$
has been reached. Its output probability differs from $\lambda/\rho$ by
less than $\theta/4$. If each implemented input race bit has conditional
TV error at most $\beta=\theta/(8n)$, the full output error is less than
$\theta/2$.
\end{lemma}
\begin{proof}
Independence of the two clocks gives
$\Prb(T<E)=\int_0^\infty\lambda e^{-(\lambda+\rho)t}\,dt$.
The infinite-walk hitting probability is the classical biased-walk solution
$r=p/(1-p)$, obtained by solving the bounded harmonic hitting equation or
optional stopping on finite intervals and taking a limit. Since $r\le1/2$,
the paths lost by absorption at $-m$ have probability at most the eventual
success probability starting there, namely $r^{m+1}\le2^{-m-1}<\theta/16$.

For the unlimited walk, $\E S_n\le-n/3$. A path not yet absorbed at either
boundary has $S_n>-m$. Since $n\ge6m$, Hoeffding for independent increments
in $[-1,1]$ gives
\begin{align*}
 \Prb(\hbox{not absorbed at }n)
 &\le\Prb(S_n-\E S_n>n/3-m)\\
 &\le\exp[-(n/3-m)^2/(2n)]\\
 &\le e^{-n/72}=e^{-m}\le2^{-m}\le\theta/8 .
\end{align*}
The reference walk uses independent ideal coins. Conditional kernel errors
of the physical implementation are coupled sequentially: the probability
any of at most $n$ races differs is at most $n\beta=\theta/8$.
Adding the three errors proves the result without assuming the physical
coins independent.
\end{proof}

The construction is a particular Bernoulli factory; the general analytic
Bernoulli-factory theory is established \cite{nacuperes}. The explicit
bounded walk, censored physical race and total error/cost accounting are
provided because they are needed by the claimed physical operation.

\subsection{A finite deadline and comparison device}

For each ideal race input at accuracy $\beta$, a sufficient finite recipe is
\begin{equation}\label{eq:timing}
 D=\rho^{-1}\lceil\log_2(32/\beta)\rceil,\qquad
 0<\Delta\le\beta/(256\rho).
\end{equation}
The ideal deadline tail is $e^{-\rho D}\le\beta/32$.
Use an independent uniform $b$-bit grid variable to implement the inverse
exponential map. Couple it to a true $U\in(0,1)$ by a uniform grid interval.
On $E=-\log U/\rho\le D$ the map's derivative is at most
$e^{\rho D}/\rho$. A sufficient uniform grid mesh is
\[
 2^{-b}\le\min\{\beta/256,\ \rho\Delta e^{-\rho D}/32\}.
\]
An endpoint cell intersecting $e^{-\rho D}$ has probability at most
$2^{-b}$ and is charged separately. The endpoint $U=0$ is never evaluated:
all grid cells below the cap threshold return the declared fallback.
On the other cells, finite interval evaluation and rounding produce a
scheduled deadline within $\Delta$ of the coupled ideal $E$. Choosing
interval width at most $\Delta/4$ and rounding to a clock grid of spacing
at most $\Delta/4$ supplies this bound. The needed bit length is
$O(\log(1/\beta)+\log(1/(\rho\Delta)))$; the displayed inequalities can be
checked with rational conservative exponential bounds.

Include actual deadline service error and timestamp/response errors in a
combined comparison error of at most $3\Delta$. The physical comparison can
then differ from the ideal event only if $|E-T|\le3\Delta$, or on a cap or
numerical failure. Conditional on $T$, the density of the independent $E$
is at most $\rho$, so the tie band has probability at most $6\rho\Delta$.
Use a fixed rule in ties. Thus cap, threshold cell, rounding and tie-band
errors sum to less than $\beta/8$ with the choices above; the remaining
budget up to $\beta$ covers declared read, missed-exit and service errors.
If these errors are deterministic accuracy bounds rather than random
failures, their contributions are just the tie-band bound. No total
variation approximation of a continuous real clock by a discrete variable
is asserted; only the event-comparison bit is approximated.

Each race ends by $D+O(\Delta)$, including a no-exit observation. It never
waits for an uncontrolled rare full dwell $1/\lambda$. But it uses a new
requested-state pin and release before every trial, even when that state
is rarely occupied naturally. The pin is a substantive resource.

\subsection{Full operating kernel and cost}

Choose a coordinate uniformly from $N$, request its actual operating
environment and neighborhood, and use Lemma~\ref{lem:walk} to flip it with
probability $\lambda_i/\rho$. The ideal kernel is exactly
$P=I+Q^u/(N\rho)$. A finite coordinate selector can use an error budget
$\theta/4$; the preceding factory has slack for the remaining budget.

For the actual source and target states,
$\pi^u(x)Q^u(x,x^{(i)})=\pi^u(x^{(i)})Q^u(x^{(i)},x)$, because the
same $r_i^u(s)$ multiplies the complementary logistic probabilities.
Therefore $P^*=P$. This route acquires both update interfaces only for the
declared reversible form. It does not construct the reversed kernel of an
arbitrary unknown nonreversible process.

A worst-case service cost for an operating update at error $\theta$ is
\begin{equation}\label{eq:flipcost}
\begin{split}
 C_P(\theta)\le{}&C_{\rm context}+C_{\rm field}(\xi_{\rm op})
                  +C_{\rm slew}+C_{\rm arith}\\
 &+n\big[C_{\rm pin}+C_{\rm settle}
       +C_{\rm event}(\beta)+C_{\rm timer}(\Delta)
       +D+O(\Delta)\big],
\end{split}
\end{equation}
with $m,n,\beta,D,\Delta$ as above and with any operating-environment
reproduction precision $\xi_{\rm op}$ required to produce the declared
actual $Q^u$. When faithful copies reproduce that environment exactly by
assumption, $C_{\rm field}(\xi_{\rm op})$ denotes that paid reproducible
setting. Otherwise the induced kernel difference is separately bounded
and allocated to $\theta$; an unknown arbitrary dependence of $r_i^u$
on the field does not permit a continuity inference from field accuracy.
Deployment must reproduce the same environment or carry an additional
proved law-error bound.

The number of pins is $O(\log(1/\theta))$, and the total censored dwell is
$O(\rho^{-1}\log^2(1/\theta))$ with the explicit conservative constants
above. Small rates need not be estimated: a small success probability is
implemented by a correspondingly unlikely success. This does not improve
the cost of \emph{estimating} a small rate relatively. It does not remove
the numerical $r_*^{-1}$ dependence of $g^{-1}$ or the cost of physically
setting a locally unlikely starting bit.
